\chapter*{Future Scope}
While the QRAG framework demonstrates promising results, several opportunities exist for further improvement and expansion:

\begin{itemize}
    \item \textbf{Scalability to Large Datasets:} Future work can extend the system to handle larger and more diverse datasets, improving generalization across domains.

    \item \textbf{Improved Quantum Circuit Design:} Advanced ansatz designs and improved feature encoding strategies can enhance the performance of quantum models.

    \item \textbf{Error Mitigation Techniques:} Incorporating noise reduction and error correction methods can improve the reliability of quantum computations on NISQ devices.

    \item \textbf{Real-Time Integration:} Optimizing the switching mechanism between classical and quantum pipelines can reduce latency and enable real-time applications.

    \item \textbf{Enhanced Evaluation Metrics:} Developing more robust evaluation metrics that capture semantic correctness beyond vector similarity can provide better performance assessment.

    \item \textbf{Broader Applications:} The hybrid framework can be extended to other domains such as recommendation systems, question answering, and scientific knowledge discovery.
\end{itemize}

\bigskip\noindent
The QRAG framework represents a step toward practical hybrid quantum-classical systems in NLP. As quantum hardware continues to evolve, such architectures have the potential to play a significant role in advancing semantic understanding and intelligent information retrieval systems.
